% GENERATED FILE. Edit canonical lessons, then run npm run generate:books.
% canonical-source-hash: fnv1a64:954122616b378e9d
% canonical-lessons: BN-C183-chulkano, BN-C183-hai-tola, BN-C183-nak-daka, BN-C183-jaga, BN-C183-snan-kora

\chapter{To scratch, To yawn, To snore, To wake up, To bathe}
\label{ch:bn-to-scratch-to-yawn-to-snore-to-wake-up-to-bathe}
\hlchaptermodality{183}

\begin{chapteropening}
\textbf{By the end of this chapter:} \emph{I can say to scratch, to yawn, to snore, to wake up and to bathe.}

Complete the last lesson of chapter 183: I can say to scratch, to yawn, to snore, to wake up and to bathe.
\end{chapteropening}

\section[\texorpdfstring{\bn{চুলকানো} (chulkāno)}{চুলকানো (chulkāno)}]{\bn{চুলকানো} (chulkāno) — to scratch}

\label{lesson:BN-C183-chulkano}



\begin{quote}
Before the new one: say the Bengali for to share, then the Bengali for to add.
\end{quote}

\subsection*{You'll want to know: \bn{চুলকানো}}
\textbf{\bn{চুলকানো}} — \emph{chulkāno} — \textquotedblleft{}to scratch\textquotedblright{}.

\begin{cousinweb}
CDIAL has no row for \textbf{\bn{চুলকানো}} (\emph{chulkāno}), but it appears to belong with \emph{*cul}, \textquotedblleft{}fidget, itch\textquotedblright{}, a form CDIAL reconstructs; under it CDIAL lists Hindi \emph{cul}, \textquotedblleft{}itching\textquotedblright{}, and \emph{culculānā}, \textquotedblleft{}to itch\textquotedblright{}, and Bengali \emph{culculiyā}, \textquotedblleft{}restless\textquotedblright{} (CDIAL 4874).
\end{cousinweb}

\begin{grammarlens}[title={What you've built}]
One more word for this chapter.
\end{grammarlens}

\subsection*{Your turn}
\begin{itemize}
\raggedright
  \item \emph{Say it:} \emph{chulkāno}
  \item \emph{Say it:} \emph{chulkāno}, once more
  \item \emph{Say it:} say \emph{chulkāno}
  \item \emph{Recall:} say the Bengali for to share, then the Bengali for to add, then say \emph{chulkāno} again
\end{itemize}

\begin{tcolorbox}[breakable,colback=teal!4,colframe=teal!35!black,title={Before you move on}]
What does \bn{চুলকানো} mean? (\textquotedblleft{}To scratch\textquotedblright{}.) Say it once more.
\end{tcolorbox}
\section[\texorpdfstring{\bn{হাই} \bn{তোলা} (hāi tolā)}{হাই তোলা (hāi tolā)}]{\bn{হাই} \bn{তোলা} (hāi tolā) — to yawn}

\label{lesson:BN-C183-hai-tola}



\begin{quote}
Before the new one: say the Bengali for to add, then the Bengali for to scratch.
\end{quote}

\subsection*{You'll want to know: \bn{হাই} \bn{তোলা}}
\textbf{\bn{হাই} \bn{তোলা}} — \emph{hāi tolā} — \textquotedblleft{}to yawn\textquotedblright{}.

\begin{cousinweb}
\textbf{\bn{হাই} \bn{তোলা}} (\emph{hāi tolā}) is \emph{hāi}, \textquotedblleft{}a yawn\textquotedblright{}, whose origin is not settled, as CDIAL has no entry for it, + \emph{tolā}, \textquotedblleft{}to raise\textquotedblright{}, from Sanskrit \emph{tōlayati}, \textquotedblleft{}lifts, weighs\textquotedblright{} (CDIAL 5979).
\end{cousinweb}

\begin{grammarlens}[title={What you've built}]
One more word for this chapter.
\end{grammarlens}

\subsection*{Your turn}
\begin{itemize}
\raggedright
  \item \emph{Say it:} \emph{hāi tolā}
  \item \emph{Say it:} \emph{hāi tolā}, once more
  \item \emph{Say it:} say \emph{hāi tolā}
  \item \emph{Recall:} say the Bengali for to add, then the Bengali for to scratch, then say \emph{hāi tolā} again
\end{itemize}

\begin{tcolorbox}[breakable,colback=teal!4,colframe=teal!35!black,title={Before you move on}]
What does \bn{হাই} \bn{তোলা} mean? (\textquotedblleft{}To yawn\textquotedblright{}.) Say it once more.
\end{tcolorbox}
\section[\texorpdfstring{\bn{নাক} \bn{ডাকা} (nāk ḍākā)}{নাক ডাকা (nāk ḍākā)}]{\bn{নাক} \bn{ডাকা} (nāk ḍākā) — to snore}

\label{lesson:BN-C183-nak-daka}



\begin{quote}
Before the new one: say the Bengali for to scratch, then the Bengali for to yawn.
\end{quote}

\subsection*{You'll want to know: \bn{নাক} \bn{ডাকা}}
\textbf{\bn{নাক} \bn{ডাকা}} — \emph{nāk ḍākā} — \textquotedblleft{}to snore\textquotedblright{}.

\begin{cousinweb}
\textbf{\bn{নাক} \bn{ডাকা}} (\emph{nāk ḍākā}) is \emph{nāk}, \textquotedblleft{}nose\textquotedblright{}, from \emph{*nakka}, \textquotedblleft{}nose\textquotedblright{} (CDIAL 6909), + \emph{ḍākā}, \textquotedblleft{}to call\textquotedblright{}, from \emph{*ḍakk}, \textquotedblleft{}shout\textquotedblright{} (CDIAL 5516), both forms CDIAL reconstructs: \textquotedblleft{}the nose calls\textquotedblright{}.
\end{cousinweb}

\begin{grammarlens}[title={What you've built}]
One more word for this chapter.
\end{grammarlens}

\subsection*{Your turn}
\begin{itemize}
\raggedright
  \item \emph{Say it:} \emph{nāk ḍākā}
  \item \emph{Say it:} \emph{nāk ḍākā}, once more
  \item \emph{Say it:} say \emph{nāk ḍākā}
  \item \emph{Recall:} say the Bengali for to scratch, then the Bengali for to yawn, then say \emph{nāk ḍākā} again
\end{itemize}

\begin{tcolorbox}[breakable,colback=teal!4,colframe=teal!35!black,title={Before you move on}]
What does \bn{নাক} \bn{ডাকা} mean? (\textquotedblleft{}To snore\textquotedblright{}.) Say it once more.
\end{tcolorbox}
\section[\texorpdfstring{\bn{জাগা} (jāgā)}{জাগা (jāgā)}]{\bn{জাগা} (jāgā) — to wake up}

\label{lesson:BN-C183-jaga}



\begin{quote}
Before the new one: say the Bengali for to yawn, then the Bengali for to snore.
\end{quote}

\subsection*{You'll want to know: \bn{জাগা}}
\textbf{\bn{জাগা}} — \emph{jāgā} — \textquotedblleft{}to wake up\textquotedblright{}.

\begin{cousinweb}
\textbf{\bn{জাগা}} (\emph{jāgā}) is a tadbhava, a word inherited from Sanskrit and reshaped by sound change. It comes from \emph{jāgrati}, \textquotedblleft{}are awake\textquotedblright{} (CDIAL 5175). CDIAL lists Bengali \emph{jāgā}, \textquotedblleft{}to wake, watch\textquotedblright{}.
\end{cousinweb}

\begin{grammarlens}[title={What you've built}]
One more word for this chapter.
\end{grammarlens}

\subsection*{Your turn}
\begin{itemize}
\raggedright
  \item \emph{Say it:} \emph{jāgā}
  \item \emph{Say it:} \emph{jāgā}, once more
  \item \emph{Say it:} say \emph{jāgā}
  \item \emph{Recall:} say the Bengali for to yawn, then the Bengali for to snore, then say \emph{jāgā} again
\end{itemize}

\begin{tcolorbox}[breakable,colback=teal!4,colframe=teal!35!black,title={Before you move on}]
What does \bn{জাগা} mean? (\textquotedblleft{}To wake up\textquotedblright{}.) Say it once more.
\end{tcolorbox}
\section[\texorpdfstring{\bn{স্নান} \bn{করা} (snān kôrā)}{স্নান করা (snān kôrā)}]{\bn{স্নান} \bn{করা} (snān kôrā) — to bathe}

\label{lesson:BN-C183-snan-kora}



\begin{quote}
Before the new one: say the Bengali for to snore, then the Bengali for to wake up.
\end{quote}

\subsection*{You'll want to know: \bn{স্নান} \bn{করা}}
\textbf{\bn{স্নান} \bn{করা}} — \emph{snān kôrā} — \textquotedblleft{}to bathe\textquotedblright{}.

\begin{cousinweb}
\textbf{\bn{স্নান} \bn{করা}} (\emph{snān kôrā}) is \emph{snān}, \textquotedblleft{}a bath\textquotedblright{}, Sanskrit \emph{snāna} taken over whole, + \emph{kôrā}, \textquotedblleft{}to do\textquotedblright{}, from Sanskrit \emph{karoti}, \textquotedblleft{}does\textquotedblright{} (CDIAL 2814). The verb inherited from the same root is \emph{nāoyā}, \textquotedblleft{}to bathe\textquotedblright{}, which CDIAL lists under Sanskrit \emph{snāti}, \textquotedblleft{}bathes\textquotedblright{} (CDIAL 13786).
\end{cousinweb}

\begin{grammarlens}[title={What you've built}]
One more word for this chapter.
\end{grammarlens}

\subsection*{Your turn}
\begin{itemize}
\raggedright
  \item \emph{Say it:} \emph{snān kôrā}
  \item \emph{Say it:} \emph{snān kôrā}, once more
  \item \emph{Say it:} say \emph{snān kôrā}
  \item \emph{Recall:} say the Bengali for to snore, then the Bengali for to wake up, then say \emph{snān kôrā} again
\end{itemize}

\begin{tcolorbox}[breakable,colback=teal!4,colframe=teal!35!black,title={Before you move on}]
What does \bn{স্নান} \bn{করা} mean? (\textquotedblleft{}To bathe\textquotedblright{}.) Say it once more.
\end{tcolorbox}
